% Options for packages loaded elsewhere
\PassOptionsToPackage{unicode}{hyperref}
\PassOptionsToPackage{hyphens}{url}
\documentclass[
]{article}
\usepackage{xcolor}
\usepackage[margin=1in]{geometry}
\usepackage{amsmath,amssymb}
\setcounter{secnumdepth}{5}
\usepackage{iftex}
\ifPDFTeX
  \usepackage[T1]{fontenc}
  \usepackage[utf8]{inputenc}
  \usepackage{textcomp} % provide euro and other symbols
\else % if luatex or xetex
  \usepackage{unicode-math} % this also loads fontspec
  \defaultfontfeatures{Scale=MatchLowercase}
  \defaultfontfeatures[\rmfamily]{Ligatures=TeX,Scale=1}
\fi
\usepackage{lmodern}
\ifPDFTeX\else
  % xetex/luatex font selection
\fi
% Use upquote if available, for straight quotes in verbatim environments
\IfFileExists{upquote.sty}{\usepackage{upquote}}{}
\IfFileExists{microtype.sty}{% use microtype if available
  \usepackage[]{microtype}
  \UseMicrotypeSet[protrusion]{basicmath} % disable protrusion for tt fonts
}{}
\makeatletter
\@ifundefined{KOMAClassName}{% if non-KOMA class
  \IfFileExists{parskip.sty}{%
    \usepackage{parskip}
  }{% else
    \setlength{\parindent}{0pt}
    \setlength{\parskip}{6pt plus 2pt minus 1pt}}
}{% if KOMA class
  \KOMAoptions{parskip=half}}
\makeatother
\usepackage{color}
\usepackage{fancyvrb}
\newcommand{\VerbBar}{|}
\newcommand{\VERB}{\Verb[commandchars=\\\{\}]}
\DefineVerbatimEnvironment{Highlighting}{Verbatim}{commandchars=\\\{\}}
% Add ',fontsize=\small' for more characters per line
\usepackage{framed}
\definecolor{shadecolor}{RGB}{248,248,248}
\newenvironment{Shaded}{\begin{snugshade}}{\end{snugshade}}
\newcommand{\AlertTok}[1]{\textcolor[rgb]{0.94,0.16,0.16}{#1}}
\newcommand{\AnnotationTok}[1]{\textcolor[rgb]{0.56,0.35,0.01}{\textbf{\textit{#1}}}}
\newcommand{\AttributeTok}[1]{\textcolor[rgb]{0.13,0.29,0.53}{#1}}
\newcommand{\BaseNTok}[1]{\textcolor[rgb]{0.00,0.00,0.81}{#1}}
\newcommand{\BuiltInTok}[1]{#1}
\newcommand{\CharTok}[1]{\textcolor[rgb]{0.31,0.60,0.02}{#1}}
\newcommand{\CommentTok}[1]{\textcolor[rgb]{0.56,0.35,0.01}{\textit{#1}}}
\newcommand{\CommentVarTok}[1]{\textcolor[rgb]{0.56,0.35,0.01}{\textbf{\textit{#1}}}}
\newcommand{\ConstantTok}[1]{\textcolor[rgb]{0.56,0.35,0.01}{#1}}
\newcommand{\ControlFlowTok}[1]{\textcolor[rgb]{0.13,0.29,0.53}{\textbf{#1}}}
\newcommand{\DataTypeTok}[1]{\textcolor[rgb]{0.13,0.29,0.53}{#1}}
\newcommand{\DecValTok}[1]{\textcolor[rgb]{0.00,0.00,0.81}{#1}}
\newcommand{\DocumentationTok}[1]{\textcolor[rgb]{0.56,0.35,0.01}{\textbf{\textit{#1}}}}
\newcommand{\ErrorTok}[1]{\textcolor[rgb]{0.64,0.00,0.00}{\textbf{#1}}}
\newcommand{\ExtensionTok}[1]{#1}
\newcommand{\FloatTok}[1]{\textcolor[rgb]{0.00,0.00,0.81}{#1}}
\newcommand{\FunctionTok}[1]{\textcolor[rgb]{0.13,0.29,0.53}{\textbf{#1}}}
\newcommand{\ImportTok}[1]{#1}
\newcommand{\InformationTok}[1]{\textcolor[rgb]{0.56,0.35,0.01}{\textbf{\textit{#1}}}}
\newcommand{\KeywordTok}[1]{\textcolor[rgb]{0.13,0.29,0.53}{\textbf{#1}}}
\newcommand{\NormalTok}[1]{#1}
\newcommand{\OperatorTok}[1]{\textcolor[rgb]{0.81,0.36,0.00}{\textbf{#1}}}
\newcommand{\OtherTok}[1]{\textcolor[rgb]{0.56,0.35,0.01}{#1}}
\newcommand{\PreprocessorTok}[1]{\textcolor[rgb]{0.56,0.35,0.01}{\textit{#1}}}
\newcommand{\RegionMarkerTok}[1]{#1}
\newcommand{\SpecialCharTok}[1]{\textcolor[rgb]{0.81,0.36,0.00}{\textbf{#1}}}
\newcommand{\SpecialStringTok}[1]{\textcolor[rgb]{0.31,0.60,0.02}{#1}}
\newcommand{\StringTok}[1]{\textcolor[rgb]{0.31,0.60,0.02}{#1}}
\newcommand{\VariableTok}[1]{\textcolor[rgb]{0.00,0.00,0.00}{#1}}
\newcommand{\VerbatimStringTok}[1]{\textcolor[rgb]{0.31,0.60,0.02}{#1}}
\newcommand{\WarningTok}[1]{\textcolor[rgb]{0.56,0.35,0.01}{\textbf{\textit{#1}}}}
\usepackage{longtable,booktabs,array}
\newcounter{none} % for unnumbered tables
\usepackage{calc} % for calculating minipage widths
% Correct order of tables after \paragraph or \subparagraph
\usepackage{etoolbox}
\makeatletter
\patchcmd\longtable{\par}{\if@noskipsec\mbox{}\fi\par}{}{}
\makeatother
% Allow footnotes in longtable head/foot
\IfFileExists{footnotehyper.sty}{\usepackage{footnotehyper}}{\usepackage{footnote}}
\makesavenoteenv{longtable}
\usepackage{graphicx}
\makeatletter
\newsavebox\pandoc@box
\newcommand*\pandocbounded[1]{% scales image to fit in text height/width
  \sbox\pandoc@box{#1}%
  \Gscale@div\@tempa{\textheight}{\dimexpr\ht\pandoc@box+\dp\pandoc@box\relax}%
  \Gscale@div\@tempb{\linewidth}{\wd\pandoc@box}%
  \ifdim\@tempb\p@<\@tempa\p@\let\@tempa\@tempb\fi% select the smaller of both
  \ifdim\@tempa\p@<\p@\scalebox{\@tempa}{\usebox\pandoc@box}%
  \else\usebox{\pandoc@box}%
  \fi%
}
% Set default figure placement to htbp
\def\fps@figure{htbp}
\makeatother
\setlength{\emergencystretch}{3em} % prevent overfull lines
\providecommand{\tightlist}{%
  \setlength{\itemsep}{0pt}\setlength{\parskip}{0pt}}
\usepackage{bookmark}
\IfFileExists{xurl.sty}{\usepackage{xurl}}{} % add URL line breaks if available
\urlstyle{same}
\hypersetup{
  pdftitle={Chapter 13},
  pdfauthor={by Lorraine Gaudio},
  hidelinks,
  pdfcreator={LaTeX via pandoc}}

\title{Chapter 13}
\usepackage{etoolbox}
\makeatletter
\providecommand{\subtitle}[1]{% add subtitle to \maketitle
  \apptocmd{\@title}{\par {\large #1 \par}}{}{}
}
\makeatother
\subtitle{ Grouped Summaries}
\author{by Lorraine Gaudio}
\date{Version June 11, 2026}

\begin{document}
\maketitle

{
\setcounter{tocdepth}{2}
\tableofcontents
}
\section{Overview}\label{overview}

In Chapters 11 and 12, you learned how to subset and transform data with
\texttt{dplyr} using functions such as \texttt{select()},
\texttt{filter()}, \texttt{mutate()}, and \texttt{case\_when()}. In this
chapter, you will extend that workflow by learning how to organize data
into groups and then calculate summaries for those groups. You will use
\texttt{count()} to make frequency tables, \texttt{summarize()} to
compute statistics such as mean and standard deviation, and
\texttt{group\_by()} to tell R which subgroups to analyze separately.
You will also practice computing multiple statistics at once and working
with more than one grouping variable.

This chapter is about learning when and how to ask R for grouped
summaries, and then explaining what those summaries mean. You will also
practice curiosity as a learning skill. Curiosity helps you pause before
running code, predict what the output might show, check whether the
result matches your question, and decide when a side question is worth
exploring or when it is better to stay focused on the chapter goal.

At the same time, you will continue building reproducible workflow
habits in Quarto by using headings, rendering often, checking your
output, and memoing what your code is doing and what you learned from
it.

By the end of Chapter 13 you will be able to:

\begin{itemize}
\item
  Use \texttt{count()} to create frequency tables for one categorical
  variable.
\item
  Use \texttt{count()} with more than one variable to count combinations
  of categories.
\item
  Use \texttt{summarize()} to calculate summary statistics such as mean,
  standard deviation, maximum, and row count.
\item
  Use \texttt{group\_by()} to compute separate results for each group
  instead of one result for the whole dataset.
\item
  Explain how grouping changes the meaning of a summary result.
\item
  Combine \texttt{group\_by()} and \texttt{summarize()} in a pipeline to
  calculate grouped summaries.
\item
  Compute multiple summary statistics in a single \texttt{summarize()}
  call.
\item
  Distinguish between a grouping variable and a summary variable when
  analyzing data.
\item
  Choose when \texttt{count()} is the better tool and when
  \texttt{summarize(n\ =\ n())} is the better tool.
\item
  Remove grouping after summarizing by using \texttt{.groups\ =\ "drop"}
  or \texttt{ungroup()}.
\item
  Practice curiosity by making predictions, checking output, asking
  useful follow-up questions, and deciding when to stay focused.
\item
  Continue using Quarto to organize code, output, checks, and written
  interpretation in a reproducible analysis document.
\end{itemize}

Keep these goals in mind as you move through each section.

\section{Setting Up Your Notes}\label{setting-up-your-notes}

In Chapter 11, you learned how to create a Quarto file, keep the YAML
header, use Markdown headings, write memos outside code chunks, and
control chunks with \texttt{\#\textbar{}} options. In Chapter 12, you
added a table of contents, section numbering, and inline R code inside
your sentences. In Chapter 13, you will keep practicing those Quarto
habits while learning how to count categories and summarize data by
groups.

\subsection{Create a Quarto File}\label{create-a-quarto-file}

For those of you who \textbf{have GitHub} and a course project, navigate
to your course folder and open your course project. This will open
RStudio. Click ``Pull'' in the Git pane to get the latest version of the
course materials. After completing work, save, stage, commit, and push.

If you \textbf{don't have GitHub} or a course project, just open RStudio
and set your working directory to your course folder.

Open a new Quarto document in your course folder and save it as:
\texttt{chapter13\_notes.qmd}
(\texttt{File\ \textgreater{}\ New\ File\ \textgreater{}\ Quarto\ Document}).

\subsection{YAML Header}\label{yaml-header}

Copy and paste the YAML header from this file into your new Quarto
document. Update the title, author, and date fields in the YAML header.

\begin{Shaded}
\begin{Highlighting}[]
\PreprocessorTok{{-}{-}{-}}
\FunctionTok{title}\KeywordTok{:}\AttributeTok{ }\StringTok{"Chapter 13 Notes"}
\FunctionTok{subtitle}\KeywordTok{:}\AttributeTok{ }\StringTok{"group\_by() and summarize()"}
\FunctionTok{author}\KeywordTok{:}\AttributeTok{ }\StringTok{"YOUR NAME"}
\FunctionTok{date}\KeywordTok{:}\AttributeTok{ today}
\FunctionTok{format}\KeywordTok{:}
\AttributeTok{  }\FunctionTok{html}\KeywordTok{:}
\AttributeTok{    }\FunctionTok{toc}\KeywordTok{:}\AttributeTok{ }\CharTok{true}
\AttributeTok{    }\FunctionTok{toc{-}depth}\KeywordTok{:}\AttributeTok{ }\DecValTok{2}
\AttributeTok{    }\FunctionTok{number{-}sections}\KeywordTok{:}\AttributeTok{ }\CharTok{true}
\AttributeTok{    }\FunctionTok{embed{-}resources}\KeywordTok{:}\AttributeTok{ }\CharTok{true}
\PreprocessorTok{{-}{-}{-}}
\end{Highlighting}
\end{Shaded}

\subsection{Headings}\label{headings}

Use the section headings in this file to create the structure of your
notes. You can copy and paste the section headings into your new Quarto
document.

\begin{Shaded}
\begin{Highlighting}[]

\FunctionTok{\# Loading Packages}

\FunctionTok{\# Counts, Center, and Spread}

\FunctionTok{\# count()}

\FunctionTok{\# summarize()}

\FunctionTok{\# group\_by()}

\FunctionTok{\# Multiple Statistics}

\FunctionTok{\# Multiple Variables}
\end{Highlighting}
\end{Shaded}

Work through the sections under each heading. \textbf{Tip:} Render to
HTML often to see the results.

\subsection{\texorpdfstring{ Packages}{ Packages}}\label{packages}

Install once (if needed): \texttt{install.packages("dplyr")};
\texttt{install.packages("dslabs")}

Load the packages

\begin{Shaded}
\begin{Highlighting}[]
\CommentTok{\#  Type this in a Quarto code chunk:}
\FunctionTok{library}\NormalTok{(dslabs)  }\CommentTok{\# Data science labs package}
\FunctionTok{library}\NormalTok{(dplyr)   }\CommentTok{\# Data manipulation package}
\end{Highlighting}
\end{Shaded}

\section{\texorpdfstring{ Counts, Center, and
Spread}{ Counts, Center, and Spread}}\label{counts-center-and-spread}

Before using \texttt{summarize()} and \texttt{count()}, it helps to
review a few data words you have likely seen before. You do \textbf{not}
need to calculate these by hand in this chapter. You do need to
understand what they mean so you can choose the right function and
explain your results.

Curiosity as an intellectual virtue means choosing where to spend your
attention. If these review terms feel rusty, pause here and use this
section to refresh them. If you already practiced these ideas (e.g., in
the Data Understanding Notebook activity), you can skip this review
section and \hyperref[dplyr-count]{jump to dplyr's \texttt{count()}}.
Return if you need a reminder.

\subsection{Data Terminology}\label{data-terminology}

A \textbf{numeric variable} stores numbers that can be measured or
calculated, such as \texttt{mpg}, height, or test score.

A \textbf{categorical variable} stores labels or group names, such as
state, disease, or cylinder count.

In this chapter, we often \textbf{group by a categorical variable} and
then \textbf{summarize a numeric variable}.

\subsection{Count}\label{count}

A \textbf{count} tells you how many observations there are.

\begin{itemize}
\item
  For a whole dataset, count means the total number of rows.
\item
  For grouped data, count means the number of rows in each group.
\end{itemize}

If you want to know how many cars have 4 cylinders, 6 cylinders, or 8
cylinders, you are asking for a \textbf{count}.

\subsection{Center}\label{center}

\textbf{Center} describes a typical value.

\subsubsection{Mean}\label{mean}

The \textbf{mean} is the average.

It adds the values and divides by how many values there are.

If the mean \texttt{mpg} for cars with 4 cylinders is higher than the
mean \texttt{mpg} for cars with 8 cylinders, then the 4-cylinder cars
are more fuel efficient on average.

\subsubsection{Median}\label{median}

The \textbf{median} is the middle value when the data are sorted.

Median is useful when extreme values might pull the mean too high or too
low.

\subsection{Spread}\label{spread}

\textbf{Spread} describes how much the values vary.

\subsubsection{Standard Deviation}\label{standard-deviation}

The \textbf{standard deviation} tells you how far values tend to be from
the mean.

\begin{itemize}
\item
  A \textbf{small} standard deviation means the values are packed
  closely together.
\item
  A \textbf{large} standard deviation means the values are more spread
  out.
\end{itemize}

If one cylinder group has a much larger standard deviation for
\texttt{mpg}, that group has more variability in fuel efficiency.

\subsubsection{Range}\label{range}

The \textbf{range} is the distance from the smallest value to the
largest value.

Range gives a quick sense of spread, but it only uses the minimum and
maximum values.

\begin{center}\rule{0.5\linewidth}{0.5pt}\end{center}

The goal is to know which summary tool matches the question.

\begin{itemize}
\item
  use \textbf{count} when you want to know how many rows are in a
  dataset or group
\item
  use \textbf{mean} when you want an average
\item
  use \textbf{standard deviation} when you want to describe spread
\item
  use a \textbf{categorical variable} to make groups
\item
  use a \textbf{numeric variable} when calculating statistics like mean
  or standard deviation
\end{itemize}

We'll start with counting categories with \texttt{count()} because
counts help you see what categories are present before you calculate
more complex summaries.

\section{dplyr's count()}\label{dplyr-count}

You can use \texttt{count()} to quickly count rows for each category or
combination of categories.

\texttt{count()} is especially useful when you want a \textbf{frequency
table}.

\texttt{?count}

The SYNTAX

\texttt{data\ \%\textgreater{}\%\ count(variable\_name)}

Or, for more than one grouping variable:

\texttt{data\ \%\textgreater{}\%\ count(variable1,\ variable2,\ ...)}

\subsection{Count one variable}\label{count-one-variable}

Count how many cars have each cylinder value.

\begin{Shaded}
\begin{Highlighting}[]
\CommentTok{\#  Type this in a Quarto code chunk:}
\NormalTok{mtcars }\SpecialCharTok{\%\textgreater{}\%}
  \FunctionTok{count}\NormalTok{(cyl)}
\end{Highlighting}
\end{Shaded}

\begin{verbatim}
##   cyl  n
## 1   4 11
## 2   6  7
## 3   8 14
\end{verbatim}

The code above counts how many rows fall into each value of
\texttt{cyl}. The result is a new table with one row per cylinder value
and a column named \texttt{n} that stores the counts.

This is the curiosity move here: before I compare or explain anything, I
want to know what categories are in the data and how many rows belong to
each one. This step helps me understand the data before I ask a bigger
data question.

\textbf{Test knit.} Does the table show one row for each cylinder value
and a column named \texttt{n}?

\subsubsection{count() versus length()}\label{count-versus-length}

\texttt{length()} counts how many values are in a vector.

\texttt{count()} counts rows in a data frame for each category or
combination of categories.

\subsubsection{count() with sorting}\label{count-with-sorting}

Sometimes you want to see the biggest groups first.

Count cars by cylinder and sort from largest count to smallest count.

\begin{Shaded}
\begin{Highlighting}[]
\CommentTok{\#  Type this in a Quarto code chunk:}
\NormalTok{mtcars }\SpecialCharTok{\%\textgreater{}\%}
  \FunctionTok{count}\NormalTok{(cyl, }\AttributeTok{sort =} \ConstantTok{TRUE}\NormalTok{)}
\end{Highlighting}
\end{Shaded}

\begin{verbatim}
##   cyl  n
## 1   8 14
## 2   4 11
## 3   6  7
\end{verbatim}

\texttt{sort\ =\ TRUE} arranges the result so the largest count appears
first. The counts are the same as before; only the row order changed.

\subsection{Count multiple variables}\label{count-multiple-variables}

You can also count combinations of categories.

This is a good time to indulge curiosity a little further. Counting one
variable told us how many cars had each cylinder value. Counting two
variables lets us ask a more specific question: how are cylinder values
distributed across transmission types?

Count how many cars fall into each combination of cylinder count and
transmission type.

\begin{quote}
\textbf{Notice:} In \texttt{mtcars}, \texttt{am} is a coded transmission
variable. A value of \texttt{0} means automatic, and a value of
\texttt{1} means manual.
\end{quote}

\begin{Shaded}
\begin{Highlighting}[]
\CommentTok{\#  Type this in a Quarto code chunk:}
\NormalTok{mtcars }\SpecialCharTok{\%\textgreater{}\%}
  \FunctionTok{count}\NormalTok{(cyl, am)}
\end{Highlighting}
\end{Shaded}

\begin{verbatim}
##   cyl am  n
## 1   4  0  3
## 2   4  1  8
## 3   6  0  4
## 4   6  1  3
## 5   8  0 12
## 6   8  1  2
\end{verbatim}

\texttt{count(cyl,\ am)} creates one row for each unique combination of
\texttt{cyl} and \texttt{am}, then counts how many rows are in each
combination.

\begin{center}\rule{0.5\linewidth}{0.5pt}\end{center}

This is useful when one grouping variable is not enough.

If our question is only ``How many rows are in each category?'',
\texttt{count()} is enough. But what if our question asks about
averages, maximums, spread, or other statistics? We'll need
\texttt{summarize()}.

\section{dplyr's summarize()}\label{dplyrs-summarize}

You can use \texttt{summarize()} to compute summary statistics like the
mean, median, etc.

\texttt{?summarize}

The SYNTAX

\texttt{data\ \%\textgreater{}\%}
\texttt{summarize(new\_column\_name\ =\ some\_operation(column\_name))}

Find the overall mean mpg for all cars.

\begin{Shaded}
\begin{Highlighting}[]
\CommentTok{\#  Type this in a Quarto code chunk:}
\NormalTok{mtcars }\SpecialCharTok{\%\textgreater{}\%} 
  \FunctionTok{summarize}\NormalTok{(}\AttributeTok{mean\_mpg =} \FunctionTok{mean}\NormalTok{(mpg))}
\end{Highlighting}
\end{Shaded}

\begin{verbatim}
##   mean_mpg
## 1 20.09062
\end{verbatim}

The code above creates a new column called \texttt{mean\_mpg} that
contains the mean of the \texttt{mpg} column from the \texttt{mtcars}
dataset.

\textbf{We name the new column using
\texttt{new\_name\ =\ calculation}.}

Find several summary statistics for mpg: mean, max, and count of cars.

\begin{Shaded}
\begin{Highlighting}[]
\CommentTok{\#  Type this in a Quarto code chunk:}
\NormalTok{mtcars }\SpecialCharTok{\%\textgreater{}\%} 
  \FunctionTok{summarize}\NormalTok{(}\AttributeTok{mean\_mpg =} \FunctionTok{mean}\NormalTok{(mpg),}
            \AttributeTok{max\_mpg =} \FunctionTok{max}\NormalTok{(mpg),}
            \AttributeTok{count =} \FunctionTok{n}\NormalTok{())}
\end{Highlighting}
\end{Shaded}

\begin{verbatim}
##   mean_mpg max_mpg count
## 1 20.09062    33.9    32
\end{verbatim}

\begin{quote}
\textbf{Notice:} Do not put pipes \texttt{\%\textgreater{}\%} inside
\texttt{summarize()}. Inside the function, separate each summary
calculation with a comma \texttt{,}.
\end{quote}

\textbf{Test knit.} Does the output show one row with three summary
columns: \texttt{mean\_mpg}, \texttt{max\_mpg}, and \texttt{count}?

\begin{center}\rule{0.5\linewidth}{0.5pt}\end{center}

\texttt{summary()} is a base R function that gives a quick overall
description of an object, such as basic statistics for each variable in
a data frame. \texttt{summarize()} from \texttt{dplyr} is used to create
your own summary results, usually after \texttt{group\_by()}, such as
the mean or count for each group.

\subsection{count() versus summarize()}\label{count-versus-summarize}

You can often get the same result in two ways.

Use group\_by() and summarize() to produce the same counts by cylinder.

\begin{Shaded}
\begin{Highlighting}[]
\CommentTok{\#  Type this in a Quarto code chunk:}
\NormalTok{mtcars }\SpecialCharTok{\%\textgreater{}\%}
  \FunctionTok{group\_by}\NormalTok{(cyl) }\SpecialCharTok{\%\textgreater{}\%}
  \FunctionTok{summarize}\NormalTok{(}\AttributeTok{n =} \FunctionTok{n}\NormalTok{())}
\end{Highlighting}
\end{Shaded}

\begin{verbatim}
## # A tibble: 3 x 2
##     cyl     n
##   <dbl> <int>
## 1     4    11
## 2     6     7
## 3     8    14
\end{verbatim}

This gives the same idea as \texttt{count(cyl)}. The difference is that
\texttt{count()} is shorter and faster to write when your main goal is
just to count rows.

Use \texttt{count()} when you mainly want frequencies. Use
\texttt{group\_by()} with \texttt{summarize()} when you want to
calculate other statistics too, such as mean and standard deviation.

\section{dplyr's group\_by()}\label{dplyrs-group_by}

\texttt{group\_by()} changes the behavior of summary functions.

The SYNTAX

\texttt{data\ \%\textgreater{}\%}
\texttt{group\_by(variable1,\ variable2,\ ...)\ \%\textgreater{}\%}
\texttt{some\_operation()}

Find the mean mpg for each cylinder count.

\begin{Shaded}
\begin{Highlighting}[]
\CommentTok{\#  Type this in a Quarto code chunk:}
\NormalTok{cyl\_means }\OtherTok{\textless{}{-}}\NormalTok{ mtcars }\SpecialCharTok{\%\textgreater{}\%}
  \FunctionTok{group\_by}\NormalTok{(cyl) }\SpecialCharTok{\%\textgreater{}\%}
  \FunctionTok{summarize}\NormalTok{(}\AttributeTok{mean\_mpg =} \FunctionTok{mean}\NormalTok{(mpg))}
\end{Highlighting}
\end{Shaded}

\texttt{group\_by(cyl)} splits the cars into groups based on cylinder
count, and \texttt{summarize()} calculates \texttt{mean(mpg)} for each
group. \texttt{mean\_mpg} is the name of the new summary column.

\begin{quote}
\textbf{Notice:} We chain \texttt{group\_by()} and \texttt{summarize()}
together using the pipe operator \texttt{\%\textgreater{}\%}. First,
\texttt{group\_by(cyl)} tells R to work separately by cylinder value.
Then, \texttt{summarize()} calculates the mean \texttt{mpg} for each
group.
\end{quote}

\begin{Shaded}
\begin{Highlighting}[]
\CommentTok{\#  }
\NormalTok{cyl\_means  }
\end{Highlighting}
\end{Shaded}

\begin{verbatim}
## # A tibble: 3 x 2
##     cyl mean_mpg
##   <dbl>    <dbl>
## 1     4     26.7
## 2     6     19.7
## 3     8     15.1
\end{verbatim}

This is the main curiosity move in Chapter 13. We are no longer asking,
``What is the overall mean?'' We're asking, ``Does the mean change
across groups?'' Grouping helps us ask a sharper question without
leaving the learning goal.

\textbf{Test knit}. Does cyl\_means print as a small table with one row
per cylinder value?

Grouping cars by \texttt{cyl} is useful because cars with different
numbers of cylinders often have different fuel efficiency, so one
overall average mpg would hide those differences. By grouping first, you
can compare the typical mpg for 4-, 6-, and 8-cylinder cars separately.

\subsection{Explicit vs.~Default
Naming}\label{explicit-vs.-default-naming}

If you do not give the summary column a name, \textbf{\texttt{dplyr}
uses the summary expression itself as the column name} by default. So
\texttt{summarize(mean(mpg))} creates a column named \texttt{mean(mpg)}
instead of a cleaner name like \texttt{mean\_mpg}.

Find the mean mpg for each cylinder count without naming the summary
column.

\begin{Shaded}
\begin{Highlighting}[]
\CommentTok{\#  Type this in a Quarto code chunk:}
\NormalTok{mtcars }\SpecialCharTok{\%\textgreater{}\%}
  \FunctionTok{group\_by}\NormalTok{(cyl) }\SpecialCharTok{\%\textgreater{}\%}
  \FunctionTok{summarize}\NormalTok{(}\FunctionTok{mean}\NormalTok{(mpg))      }
\end{Highlighting}
\end{Shaded}

\begin{verbatim}
## # A tibble: 3 x 2
##     cyl `mean(mpg)`
##   <dbl>       <dbl>
## 1     4        26.7
## 2     6        19.7
## 3     8        15.1
\end{verbatim}

The default column name is the code itself: \texttt{mean(mpg)}. Default
names can be confusing when reading or sharing results, so naming
summary columns makes tables clearer.

Notice the problem and build the habit of naming summary columns
clearly.

\begin{center}\rule{0.5\linewidth}{0.5pt}\end{center}

Now that you have seen how \texttt{group\_by()} changes one summary
calculation, you can ask for several summary statistics at the same
time.

\section{Multiple Statistics}\label{multiple-stats}

Create several summary columns in one \texttt{summarize()} call,
separated by commas.

This is useful curiosity: one statistic rarely tells the whole story. A
mean tells me the center, a standard deviation tells me about spread,
and \texttt{n()} tells me how many rows were used for each group.

\emph{Use \texttt{?n} to learn about the \texttt{n()} function.}

\begin{Shaded}
\begin{Highlighting}[]
\CommentTok{\#  Type this in a Quarto code chunk:}
\NormalTok{cyl\_stats }\OtherTok{\textless{}{-}}\NormalTok{ mtcars }\SpecialCharTok{\%\textgreater{}\%}
  \FunctionTok{group\_by}\NormalTok{(cyl) }\SpecialCharTok{\%\textgreater{}\%}
  \FunctionTok{summarize}\NormalTok{(}\AttributeTok{mean\_mpg =} \FunctionTok{mean}\NormalTok{(mpg),}
            \AttributeTok{sd\_mpg   =} \FunctionTok{sd}\NormalTok{(mpg),}
            \AttributeTok{n        =} \FunctionTok{n}\NormalTok{())}
\end{Highlighting}
\end{Shaded}

\begin{Shaded}
\begin{Highlighting}[]
\CommentTok{\#  Verify}
\NormalTok{cyl\_stats}
\end{Highlighting}
\end{Shaded}

\begin{verbatim}
## # A tibble: 3 x 4
##     cyl mean_mpg sd_mpg     n
##   <dbl>    <dbl>  <dbl> <int>
## 1     4     26.7   4.51    11
## 2     6     19.7   1.45     7
## 3     8     15.1   2.56    14
\end{verbatim}

\texttt{n()} counts how many rows are in each group created by
\texttt{group\_by()}.

It is similar to \texttt{count()} because both tell you how many
observations are in a group, but \texttt{n()} is used inside
\texttt{summarize()} when you are also calculating other statistics.

\begin{quote}
\textbf{Notice:} \texttt{length()} counts the number of elements in an
object. For a vector, that means the number of values. For a data frame,
\texttt{length()} counts columns, not rows. That is why \texttt{n()} is
the clearer tool here.
\end{quote}

\textbf{Test knit.} Does the output show one row for each cylinder value
and three summary columns: \texttt{mean\_mpg}, \texttt{sd\_mpg}, and
\texttt{n}?

\begin{center}\rule{0.5\linewidth}{0.5pt}\end{center}

\section{Multiple Variables}\label{multiple-vars}

\texttt{group\_by()} with one column creates one group for each unique
value in that column. \texttt{group\_by()} with two columns creates one
group for each unique combination of values.

This is a good moment to make a prediction before running the code.
Curiosity is pausing long enough to make a reasonable guess, then
checking that guess against the output.

How many groups will \texttt{group\_by(cyl,\ am)} create?

Predict how many groups will be created by \texttt{group\_by(cyl,\ am)}
before running the code.

\begin{Shaded}
\begin{Highlighting}[]
\CommentTok{\#  Pre{-}check}
\CommentTok{\# Recall that mtcars$cyl has three unique values:}
\FunctionTok{sort}\NormalTok{(}\FunctionTok{unique}\NormalTok{(mtcars}\SpecialCharTok{$}\NormalTok{cyl))}
\end{Highlighting}
\end{Shaded}

\begin{verbatim}
## [1] 4 6 8
\end{verbatim}

\begin{Shaded}
\begin{Highlighting}[]
\CommentTok{\# Recall that mtcars$am has two unique values:}
\FunctionTok{sort}\NormalTok{(}\FunctionTok{unique}\NormalTok{(mtcars}\SpecialCharTok{$}\NormalTok{am))}
\end{Highlighting}
\end{Shaded}

\begin{verbatim}
## [1] 0 1
\end{verbatim}

Two variables with 3 and 2 unique values can create \(3 \times 2 = 6\)
possible groups.

Find the mean mpg for each combination of cylinder count and
transmission type.

\begin{quote}
\textbf{Notice:} In \texttt{mtcars}, \texttt{am} is a coded transmission
variable. A value of \texttt{0} means automatic, and a value of
\texttt{1} means manual.
\end{quote}

\begin{Shaded}
\begin{Highlighting}[]
\CommentTok{\#  Type this in a Quarto code chunk:}
\NormalTok{mtcars }\SpecialCharTok{\%\textgreater{}\%}
  \FunctionTok{group\_by}\NormalTok{(cyl, am) }\SpecialCharTok{\%\textgreater{}\%}
  \FunctionTok{summarize}\NormalTok{(}\AttributeTok{avg\_mpg =} \FunctionTok{mean}\NormalTok{(mpg),}
            \AttributeTok{count =} \FunctionTok{n}\NormalTok{())}
\end{Highlighting}
\end{Shaded}

\begin{verbatim}
## # A tibble: 6 x 4
## # Groups:   cyl [3]
##     cyl    am avg_mpg count
##   <dbl> <dbl>   <dbl> <int>
## 1     4     0    22.9     3
## 2     4     1    28.1     8
## 3     6     0    19.1     4
## 4     6     1    20.6     3
## 5     8     0    15.0    12
## 6     8     1    15.4     2
\end{verbatim}

:

\begin{itemize}
\item
  Group 1: cyl = 4, am = 0 (4-cylinder automatic cars)
\item
  Group 2: cyl = 4, am = 1 (4-cylinder manual cars)
\item
  Group 3: cyl = 6, am = 0 (6-cylinder automatic cars)
\item
  Group 4: cyl = 6, am = 1 (6-cylinder manual cars)
\item
  Group 5: cyl = 8, am = 0 (8-cylinder automatic cars)
\item
  Group 6: cyl = 8, am = 1 (8-cylinder manual cars)
\end{itemize}

\textbf{Test knit.} Does the output show one row for each \texttt{cyl}
and \texttt{am} combination?

\begin{center}\rule{0.5\linewidth}{0.5pt}\end{center}

If you see a message about grouped output, keep going. The next section
explains what that message means.

\subsection{ungroup()}\label{ungroup}

When you summarize data that has more than one grouping variable,
\texttt{dplyr} may keep part of the grouping structure in the result.
That leftover grouping can affect later code, so it is important to
check whether your summary table is still grouped.

We are asking, ``What did R keep behind the scenes, and could that
affect my next step?''

\subsubsection{A. Default behavior: keeps earlier
groups}\label{a.-default-behavior-keeps-earlier-groups}

Show that with two grouping variables, \texttt{summarize()} keeps the
earlier group by default.

\begin{Shaded}
\begin{Highlighting}[]
\CommentTok{\#  Type this in a Quarto code chunk:}
\CommentTok{\# A) Default behavior: "drop\_last" (keeps earlier groups)}
\NormalTok{cyl\_gear\_stats }\OtherTok{\textless{}{-}}\NormalTok{ mtcars }\SpecialCharTok{\%\textgreater{}\%}
  \FunctionTok{group\_by}\NormalTok{(cyl, gear) }\SpecialCharTok{\%\textgreater{}\%}
  \FunctionTok{summarize}\NormalTok{(}\AttributeTok{mean\_mpg =} \FunctionTok{mean}\NormalTok{(mpg))   }\CommentTok{\# no .groups set}
\end{Highlighting}
\end{Shaded}

\begin{verbatim}
## `summarise()` has regrouped the output.
## i Summaries were computed grouped by cyl and gear.
## i Output is grouped by cyl.
## i Use `summarise(.groups = "drop_last")` to silence this message.
## i Use `summarise(.by = c(cyl, gear))` for per-operation grouping
##   (`?dplyr::dplyr_by`) instead.
\end{verbatim}

\begin{Shaded}
\begin{Highlighting}[]
\CommentTok{\#   Verify}
\FunctionTok{group\_vars}\NormalTok{(cyl\_gear\_stats)  }\CommentTok{\# what groups remain?}
\end{Highlighting}
\end{Shaded}

\begin{verbatim}
## [1] "cyl"
\end{verbatim}

\begin{Shaded}
\begin{Highlighting}[]
\CommentTok{\#  Verify}
\FunctionTok{is\_grouped\_df}\NormalTok{(cyl\_gear\_stats) }\CommentTok{\# expect TRUE}
\end{Highlighting}
\end{Shaded}

\begin{verbatim}
## [1] TRUE
\end{verbatim}

The default behavior is \texttt{.groups\ =\ "drop\_last"}. That means
\texttt{summarize()} drops the last grouping variable, \texttt{gear},
but keeps the earlier grouping variable, \texttt{cyl}.

You'll want to be explicit when you do not want any grouping left.

\subsubsection{\texorpdfstring{B. Make it explicit inside
\texttt{summarize()}}{B. Make it explicit inside summarize()}}\label{b.-make-it-explicit-inside-summarize}

Return a plain, ungrouped table by telling \texttt{summarize()} to drop
all groups.

\begin{Shaded}
\begin{Highlighting}[]
\CommentTok{\#  Type this in a Quarto code chunk:}
\CommentTok{\# B) Make it explicit: drop *all* groups in the result}
\NormalTok{cyl\_gear\_stats\_drop }\OtherTok{\textless{}{-}}\NormalTok{ mtcars }\SpecialCharTok{\%\textgreater{}\%}
  \FunctionTok{group\_by}\NormalTok{(cyl, gear) }\SpecialCharTok{\%\textgreater{}\%}
  \FunctionTok{summarize}\NormalTok{(}\AttributeTok{mean\_mpg =} \FunctionTok{mean}\NormalTok{(mpg), }\AttributeTok{.groups =} \StringTok{"drop"}\NormalTok{)}
\end{Highlighting}
\end{Shaded}

The \texttt{.groups} argument starts with a period and \texttt{"drop"}
must be written in quotation marks.

\begin{Shaded}
\begin{Highlighting}[]
\CommentTok{\#  Verify}
\FunctionTok{group\_vars}\NormalTok{(cyl\_gear\_stats\_drop)  }\CommentTok{\# expect character(0)}
\end{Highlighting}
\end{Shaded}

\begin{verbatim}
## character(0)
\end{verbatim}

\begin{Shaded}
\begin{Highlighting}[]
\CommentTok{\#  Verify}
\FunctionTok{is\_grouped\_df}\NormalTok{(cyl\_gear\_stats\_drop) }\CommentTok{\# expect FALSE}
\end{Highlighting}
\end{Shaded}

\begin{verbatim}
## [1] FALSE
\end{verbatim}

Using \texttt{.groups\ =\ "drop"} tells \texttt{summarize()} to remove
all grouping after the summary operation, resulting in a plain data
frame without any grouping structure.

\textbf{Test knit.} Does \texttt{group\_vars(cyl\_gear\_stats\_drop)}
return \texttt{character(0)} and does
\texttt{is\_grouped\_df(cyl\_gear\_stats\_drop)} return \texttt{FALSE}?

\subsubsection{C. Alternative: drop groups after
summarizing}\label{c.-alternative-drop-groups-after-summarizing}

Achieve the same ungrouped result by using \texttt{ungroup()} after
\texttt{summarize()}

\begin{Shaded}
\begin{Highlighting}[]
\CommentTok{\#  Type this in a Quarto code chunk:}
\CommentTok{\# C) Alternative: drop groups after summarizing}
\NormalTok{cyl\_gear\_stats\_ungroup }\OtherTok{\textless{}{-}}\NormalTok{ mtcars }\SpecialCharTok{\%\textgreater{}\%}
  \FunctionTok{group\_by}\NormalTok{(cyl, gear) }\SpecialCharTok{\%\textgreater{}\%}
  \FunctionTok{summarize}\NormalTok{(}\AttributeTok{mean\_mpg =} \FunctionTok{mean}\NormalTok{(mpg)) }\SpecialCharTok{\%\textgreater{}\%}
  \FunctionTok{ungroup}\NormalTok{()}
\end{Highlighting}
\end{Shaded}

\begin{verbatim}
## `summarise()` has regrouped the output.
## i Summaries were computed grouped by cyl and gear.
## i Output is grouped by cyl.
## i Use `summarise(.groups = "drop_last")` to silence this message.
## i Use `summarise(.by = c(cyl, gear))` for per-operation grouping
##   (`?dplyr::dplyr_by`) instead.
\end{verbatim}

\begin{Shaded}
\begin{Highlighting}[]
\CommentTok{\#   Verify}
\FunctionTok{group\_vars}\NormalTok{(cyl\_gear\_stats\_ungroup)  }\CommentTok{\# expect character(0)}
\end{Highlighting}
\end{Shaded}

\begin{verbatim}
## character(0)
\end{verbatim}

\begin{Shaded}
\begin{Highlighting}[]
\CommentTok{\#  Verify}
\FunctionTok{is\_grouped\_df}\NormalTok{(cyl\_gear\_stats\_ungroup) }\CommentTok{\# expect FALSE}
\end{Highlighting}
\end{Shaded}

\begin{verbatim}
## [1] FALSE
\end{verbatim}

Using \texttt{ungroup()} after \texttt{summarize()} removes all grouping
from the resulting data frame. This gives the same ungrouped result as
using \texttt{.groups\ =\ "drop"} inside \texttt{summarize()}.

When we are done summarizing by groups, we want remove grouping for a
plain table.

Pick one method and use it consistently: - Use
\texttt{.groups\ =\ "drop"} when you want to remove grouping inside
\texttt{summarize()}. - Use \texttt{ungroup()} when you want to remove
grouping as a separate step after summarizing.

\begin{center}\rule{0.5\linewidth}{0.5pt}\end{center}

Speaking of summarizing, we've made it to the end of the chapter. Let's
review what you learned.

\section{Summary}\label{summary}

\section{Summary}\label{summary-1}

In this chapter, you learned how to describe data by groups using
\texttt{dplyr}. You reviewed the basic ideas of count, center, and
spread so you could interpret grouped results, then used
\texttt{count()} to make frequency tables, \texttt{summarize()} to
compute statistics, and \texttt{group\_by()} to split a dataset into
subgroups before summarizing. You also practiced calculating multiple
statistics in one call, using \texttt{n()} inside \texttt{summarize()}
to count rows within groups, and working with combinations of grouping
variables.

Most importantly, you learned that grouping changes what a summary
means. Instead of one result for the whole dataset, you can produce
separate results for each category or combination of categories. This is
a core skill for later data analysis because analysts often need to
compare subgroups, not just describe an entire dataset with one number.

You also practiced curiosity as part of your workflow. In this chapter,
curiosity meant pausing to ask what a summary is actually showing,
predicting how many groups R might create, checking whether the output
matched your expectation, and deciding when to investigate further
versus when to stay focused on the current learning goal.

\begin{center}\rule{0.5\linewidth}{0.5pt}\end{center}

\begin{center}\rule{0.5\linewidth}{0.5pt}\end{center}

Table 1. Functions for grouped summaries

{\def\LTcaptype{none} % do not increment counter
\begin{longtable}[]{@{}
  >{\raggedright\arraybackslash}p{(\linewidth - 2\tabcolsep) * \real{0.1286}}
  >{\raggedright\arraybackslash}p{(\linewidth - 2\tabcolsep) * \real{0.8714}}@{}}
\toprule\noalign{}
\begin{minipage}[b]{\linewidth}\raggedright
Function
\end{minipage} & \begin{minipage}[b]{\linewidth}\raggedright
Purpose
\end{minipage} \\
\midrule\noalign{}
\endhead
\bottomrule\noalign{}
\endlastfoot
\texttt{count()} & Counts how many rows fall into each value or
combination of values. Useful for frequency tables. \\
\texttt{group\_by()} & Creates groups based on one or more variables so
later summary functions work separately for each group. \\
\texttt{summarize()} & Creates a smaller summary table by calculating
one or more summary statistics. \\
\texttt{n()} & Counts the number of rows in the current group when used
inside \texttt{summarize()}. \\
\texttt{mean()} & Calculates the average of numeric values. \\
\texttt{sd()} & Calculates the standard deviation, which describes
spread around the mean. \\
\texttt{max()} & Finds the largest value in a vector. \\
\texttt{unique()} & Shows the distinct values in a vector. Useful for
checking possible groups before summarizing. \\
\texttt{length()} & Counts the number of elements in an object. For a
vector, it counts values. For a data frame, it counts columns, not
rows. \\
\texttt{.groups\ =\ "drop"} & Removes grouping inside
\texttt{summarize()} so the result is returned as an ungrouped table. \\
\texttt{ungroup()} & Removes grouping after a grouped operation has
already been created. \\
\end{longtable}
}

\section{Chapter Terms}\label{chapter-terms}

\textbf{Categorical Data}: Data whose values place observations into
groups or categories. The categories may have no natural order (nominal)
or may have an order with uneven gaps between categories (ordinal). In
R, categorical data are often stored as factors.

\textbf{Center}: The middle or typical location of a dataset. Measures
of center, such as the mean and median, summarize where values tend to
cluster.

\textbf{Count}: The number of observations or values in a dataset,
group, or category. In introductory statistics, count tells you how many
cases are present.

\textbf{Frequency}: The number of times a value or category appears in a
dataset.

\textbf{Mean}: The arithmetic average of a set of numbers, found by
adding all values and dividing by the number of values.

\textbf{Median}: The middle value in an ordered dataset. If there is an
even number of values, the median is the average of the two middle
values.

\textbf{Numerical Data}: Data recorded as numbers, where the values have
meaningful magnitudes and consistent intervals. Numerical data are also
called \emph{quantitative} data.

\textbf{Pipe}: An operator that passes the result of one step directly
into the next step, letting you write code as a readable sequence of
actions instead of nesting functions inside each other. In tidyverse
code, the pipe is usually \texttt{\%\textgreater{}\%}, so
\texttt{x\ \%\textgreater{}\%\ f()} means ``take \texttt{x}, then apply
\texttt{f()}''; in newer base R, a similar built-in pipe is
\texttt{\textbar{}\textgreater{}}.

\textbf{Range}: A simple measure of spread found by subtracting the
smallest value from the largest value in a dataset.

\textbf{Spread}: How much the values in a dataset vary or scatter rather
than cluster tightly together. Common measures of spread include range,
interquartile range (IQR), variance, and standard deviation.

\textbf{Standard Deviation}: A measure of spread that describes how far
values typically vary from the mean. A small standard deviation means
values are clustered close to the mean, while a large standard deviation
means they are more spread out.

\section{\texorpdfstring{ Practice
Space}{ Practice Space}}\label{practice-space}

This Practice Space helps you turn Chapter 13 concepts into a
reproducible analysis workflow.

In this chapter, you practiced using \texttt{group\_by()},
\texttt{summarize()}, \texttt{count()}, and \texttt{n()} to describe
data. You worked with grouped summaries, multiple statistics, and
combinations of categorical variables.

This work matters because analysts often need to answer questions such
as:

\begin{itemize}
\item
  How many observations are in each group?
\item
  What is the average value within each group?
\item
  How much do values vary within each group?
\item
  What happens when we group by more than one variable?
\end{itemize}

These tasks give you practice writing code, checking your output, and
explaining what your results mean in words. That is part of doing
reproducible work that someone else can follow and re-run later.

\subsection{Setup}\label{setup}

\begin{enumerate}
\def\labelenumi{\arabic{enumi}.}
\tightlist
\item
  Create a new Quarto file in your course folder and save it as:
\end{enumerate}

\begin{itemize}
\item
  \texttt{chapter13\_practice.qmd}
\item
  Test Render your document often to see how your Practice Space report
  is coming together.
\item
  Submit \texttt{chapter13\_practice.html}
\end{itemize}

\begin{enumerate}
\def\labelenumi{\arabic{enumi}.}
\setcounter{enumi}{1}
\item
  Use the heading structure taught earlier in the chapter. Your Quarto
  setup for Chapter 13 should include headings, section numbering, and a
  table of contents, and you should render often to check both
  formatting and errors.
\item
  Library it up!
\end{enumerate}

\begin{itemize}
\tightlist
\item
  Make sure there is script in your document that loads \texttt{dplyr}
  and \texttt{dslabs} packages so their functions / datasets load.
\end{itemize}

\begin{enumerate}
\def\labelenumi{\arabic{enumi}.}
\setcounter{enumi}{3}
\tightlist
\item
  Memo your work with the // /
\end{enumerate}

\begin{itemize}
\item
  a memo before the code chunk that states the goal of the code and what
  you are trying to learn or understand
\item
  the code chunk itself
\item
  any needed checks after the code
\item
  a memo after the code chunk that explains what you learned,
  discovered, or confirmed
\item
  answers to all comment questions included in the task
\end{itemize}

If you used help while working, document it in a memo. For example, if
you used an LLM or another outside resource to understand an error or
concept, briefly disclose that use and explain what help you received.
Type all code yourself. This is part of the course expectation for
AI-supported learning and memo-based workflow.

\begin{center}\rule{0.5\linewidth}{0.5pt}\end{center}

Follow this template for each coding task.

What is the goal of the code? What am I trying to learn, check, or
demonstrate?

\begin{Shaded}
\begin{Highlighting}[]
\CommentTok{\#Code }
\NormalTok{your\_code }\OtherTok{\textless{}{-}} \StringTok{"here"}
\end{Highlighting}
\end{Shaded}

\begin{Shaded}
\begin{Highlighting}[]
\CommentTok{\#  Include Checks after and pre{-}checks as needed}
\FunctionTok{colnames}\NormalTok{(your\_code)}
\FunctionTok{head}\NormalTok{(your\_code)}
\FunctionTok{table}\NormalTok{(your\_code)}
\FunctionTok{class}\NormalTok{(your\_code)}
\FunctionTok{unique}\NormalTok{(your\_code)}
\end{Highlighting}
\end{Shaded}

/ What I learned / Next step. 1--2 sentences.

\begin{center}\rule{0.5\linewidth}{0.5pt}\end{center}

\subsection{Task 1}\label{task-1}

\textbf{Single grouping}

Using \texttt{stars}, compute \texttt{mean\_temperature} \textbf{AND}
\texttt{mean\_magnitude} for each type. Name the result
\textbf{\texttt{twinkle\_twinkle}}.

\begin{Shaded}
\begin{Highlighting}[]
\FunctionTok{data}\NormalTok{(stars)}
\NormalTok{?stars}
\FunctionTok{colnames}\NormalTok{(stars)}
\end{Highlighting}
\end{Shaded}

Restate the Goal. What are you trying to learn, check, or demonstrate?

\begin{Shaded}
\begin{Highlighting}[]
\CommentTok{\# Code}
\NormalTok{twinkle\_twinkle }\OtherTok{\textless{}{-}}\NormalTok{ stars }\SpecialCharTok{\%\textgreater{}\%}
  \FunctionTok{\_\_\_\_\_}\NormalTok{(\_\_\_\_\_) }\SpecialCharTok{\%\textgreater{}\%}
  \FunctionTok{summarize}\NormalTok{(}
    \AttributeTok{mean\_temperature =} \FunctionTok{\_\_\_\_\_}\NormalTok{(\_\_\_\_\_),}
    \AttributeTok{mean\_magnitude   =} \FunctionTok{\_\_\_\_\_}\NormalTok{(\_\_\_\_\_),}
    \AttributeTok{.groups =} \StringTok{"drop"}
\NormalTok{  )}
\end{Highlighting}
\end{Shaded}

\begin{Shaded}
\begin{Highlighting}[]
\CommentTok{\#  Checks}
\NormalTok{twinkle\_twinkle}
\end{Highlighting}
\end{Shaded}

Comment on the output.

\subsection{Task 2}\label{task-2}

\textbf{Single grouping + three stats}

Pick one numeric fatty acid in \texttt{olive} such as \texttt{palmitic},
\texttt{oleic}, or another fatty-acid variable. Group by \texttt{region}
then calculate

\begin{itemize}
\item
  the mean of your chosen fatty acid,
\item
  the standard deviation of your choser fatty acid, and
\item
  the number of rows in each region.
\end{itemize}

Store as \textbf{\texttt{Olive\_Garden}}.

\begin{Shaded}
\begin{Highlighting}[]
\FunctionTok{data}\NormalTok{(olive)}
\NormalTok{?olive}
\FunctionTok{colnames}\NormalTok{(olive)}
\end{Highlighting}
\end{Shaded}

Restate the goal. What are you trying to learn, check, or demonstrate?

\begin{Shaded}
\begin{Highlighting}[]
\CommentTok{\# Code}
\NormalTok{fatty\_acid }\OtherTok{\textless{}{-}} \StringTok{"\_\_\_\_\_\_\_"} 

\NormalTok{Olive\_Garden }\OtherTok{\textless{}{-}}\NormalTok{ olive \_\_\_}
  \FunctionTok{group\_by}\NormalTok{(region) \_\_\_}
  \FunctionTok{summarize}\NormalTok{(}
    \AttributeTok{mean\_value =} \FunctionTok{mean}\NormalTok{(fatty\_acid)\_\_\_}
    \AttributeTok{sd\_value   =} \FunctionTok{sd}\NormalTok{(fatty\_acid)\_\_\_}
    \AttributeTok{n          =} \FunctionTok{n}\NormalTok{()}
\NormalTok{  \_\_\_}
\end{Highlighting}
\end{Shaded}

\begin{Shaded}
\begin{Highlighting}[]
\CommentTok{\#  Checks}
\NormalTok{Olive\_Garden}
\end{Highlighting}
\end{Shaded}

In a short comment, note the region with the highest mean value for the
fatty acid you selected.

Why is region the grouping variable and the fatty acid column the
summary variable?

\subsection{Task 3}\label{task-3}

\textbf{Counting unique values}

Which function can you use to count unique values in a column?

For \texttt{us\_contagious\_diseases} from \texttt{dslabs}, determine
the number of unique diseases in the column \texttt{disease}.

\begin{Shaded}
\begin{Highlighting}[]
\FunctionTok{data}\NormalTok{(}\StringTok{"us\_contagious\_diseases"}\NormalTok{)}
\NormalTok{?us\_contagious\_diseases}
\end{Highlighting}
\end{Shaded}

\begin{Shaded}
\begin{Highlighting}[]
\CommentTok{\#  Pre{-}check}
\DocumentationTok{\#\# What are the unique diseases}
\FunctionTok{unique}\NormalTok{(us\_contagious\_diseases}\SpecialCharTok{$}\NormalTok{disease)}

\DocumentationTok{\#\# How many unique diseases are there?}
\FunctionTok{\_\_\_\_\_}\NormalTok{(}\FunctionTok{unique}\NormalTok{(us\_contagious\_diseases}\SpecialCharTok{$}\NormalTok{disease))}
\end{Highlighting}
\end{Shaded}

Inline memo sentence

Write this as normal text in your Quarto document, not as code. Replace
the blank with the correct function to count unique diseases.

\begin{Shaded}
\begin{Highlighting}[]
\NormalTok{ There are }\StringTok{\textasciigrave{}}\AttributeTok{r \_\_\_\_\_(unique(us\_contagious\_diseases$disease))}\StringTok{\textasciigrave{}}\NormalTok{ unique diseases within the }\StringTok{\textasciigrave{}}\AttributeTok{disease}\StringTok{\textasciigrave{}}\NormalTok{ column of the }\StringTok{\textasciigrave{}}\AttributeTok{us\_contagious\_diseases}\StringTok{\textasciigrave{}}\NormalTok{ dataset. They are }\StringTok{\textasciigrave{}}\AttributeTok{r unique(us\_contagious\_diseases$disease)"}\StringTok{\textasciigrave{}}
\end{Highlighting}
\end{Shaded}

For \texttt{us\_contagious\_diseases}, determine the number of states in
the column \texttt{state}.

\begin{Shaded}
\begin{Highlighting}[]
\CommentTok{\#  Pre{-}check}
\DocumentationTok{\#\# What are the unique state}
\FunctionTok{unique}\NormalTok{(us\_contagious\_diseases}\SpecialCharTok{$}\NormalTok{state)}

\DocumentationTok{\#\# How many unique diseases are there?}
\FunctionTok{\_\_\_\_\_}\NormalTok{(}\FunctionTok{unique}\NormalTok{(us\_contagious\_diseases}\SpecialCharTok{$}\NormalTok{state))}
\end{Highlighting}
\end{Shaded}

Inline memo sentence

Write this as normal text in your Quarto document, not as code. Replace
the blank with the correct function to count unique states.

\begin{Shaded}
\begin{Highlighting}[]
\NormalTok{ There are }\StringTok{\textasciigrave{}}\AttributeTok{r \_\_\_\_\_(unique(us\_contagious\_diseases$state))}\StringTok{\textasciigrave{}}\NormalTok{ unique states within the }\StringTok{\textasciigrave{}}\AttributeTok{state}\StringTok{\textasciigrave{}}\NormalTok{ column of the }\StringTok{\textasciigrave{}}\AttributeTok{us\_contagious\_diseases}\StringTok{\textasciigrave{}}\NormalTok{ dataset. }
\end{Highlighting}
\end{Shaded}

\subsection{Task 4}\label{task-4}

\textbf{Make and test a prediction}

How many rows would you expect if each state has 6 diseases?

\subsubsection{Step 1}\label{step-1}

If there are 51 states in the dataset, and you assume every state has
data for all 7 diseases, how many groups will
\texttt{group\_by(state,\ disease)} create?

A. \(51 + 7 = 58\) groups

B. \(51 \times 7 = 357\) groups

C. \(\frac{51}{7} \approx 7\) groups

D. \(51 - 7 = 44\) groups

E. \(51^7\) groups

Explain your answer choice.

\subsubsection{Step 2}\label{step-2}

Test it.

Use \texttt{group\_by()} both \texttt{state} and \texttt{disease} and
\texttt{count()} to count the number of rows for each state and disease
combination. Store as \texttt{categorical\_counts}.

Restate the Goal. What are you trying to learn, check, or demonstrate?

\begin{Shaded}
\begin{Highlighting}[]
\CommentTok{\#  Code}
\NormalTok{categorical\_counts }\OtherTok{\textless{}{-}}\NormalTok{ us\_contagious\_diseases }\SpecialCharTok{\%\textgreater{}\%}
  \FunctionTok{\_\_\_\_\_}\NormalTok{(state, disease) }\SpecialCharTok{\%\textgreater{}\%}
  \FunctionTok{\_\_\_\_\_}\NormalTok{()}
\end{Highlighting}
\end{Shaded}

\begin{Shaded}
\begin{Highlighting}[]
\CommentTok{\#  Check row count}
\FunctionTok{\_\_\_\_\_}\NormalTok{(categorical\_counts) }
\end{Highlighting}
\end{Shaded}

Comment about the results. Was your prediction of group size correct?

Explain why \texttt{group\_by()} is useful when reporting the
\textbf{frequency} of categorical data like \texttt{state} and
\texttt{disease}.

\subsection{Task 5}\label{task-5}

\textbf{National mean counts by disease}

From \texttt{us\_contagious\_diseases}, calculate the national mean of
\texttt{count} for each \texttt{disease}. Then keep only the three
diseases with the highest national mean counts. Store as
\textbf{\texttt{Infectious\_Burden}}.

\textbf{Tip:} Arrange the diseases by their national mean and keep the
top three.

\begin{Shaded}
\begin{Highlighting}[]
\CommentTok{\# National means per disease}
\NormalTok{\_\_\_\_\_ }\OtherTok{\textless{}{-}}\NormalTok{ \_\_\_\_\_ }\SpecialCharTok{\%\textgreater{}\%}
  \FunctionTok{group\_by}\NormalTok{(\_\_\_\_\_) }\SpecialCharTok{\%\textgreater{}\%}
  \FunctionTok{summarize}\NormalTok{(}
    \AttributeTok{\_\_\_\_\_ =} \FunctionTok{mean}\NormalTok{(\_\_\_\_\_, }\AttributeTok{na.rm =} \ConstantTok{TRUE}\NormalTok{),}
    \AttributeTok{.groups =} \StringTok{"drop"}
\NormalTok{  ) }\SpecialCharTok{\%\textgreater{}\%}
  \FunctionTok{arrange}\NormalTok{(}\FunctionTok{desc}\NormalTok{(\_\_\_\_\_))}
\end{Highlighting}
\end{Shaded}

\begin{Shaded}
\begin{Highlighting}[]
\CommentTok{\#  Checks}
\FunctionTok{head}\NormalTok{(Infectious\_Burden, }\DecValTok{3}\NormalTok{)}
\end{Highlighting}
\end{Shaded}

Comment on the results. Which three diseases had the highest national
mean counts?

\subsection{Task 6}\label{task-6}

\textbf{Reflection}

Write a short paragraph explaining why \texttt{group\_by()} needs to
come before \texttt{summarize()} when you want separate results for each
group.

Your paragraph should explain the difference between:

\begin{itemize}
\item
  summarizing the whole dataset
\item
  summarizing within groups
\end{itemize}

\end{document}
